% !TeX root = ../drafts/09-recursive-aggregation.tex
\section{Recursion}\label{sec:recursion}

A node proves it verified $n_{\mathrm{rec}}$ sub-proofs and $n_{\mathrm{raw}}$ XMSS signatures against one message and epoch, and publishes the sorted deduplicated union of their signer sets. The sub-proofs are proofs of the same bytecode, so recursion is self-reference: only the bytecode's \emph{size} needs a fixed point, its digest riding the public statement instead of the code. A parent rebuilds each child's statement, pinning it to the same bytecode, message and epoch, and writes one write-once cell per declared signer, holding a running count that is totalled at the end: every declared signer is therefore backed by a signature or by a verified child, and a key a child repeats takes a duplicate slot past the set, which merges overlapping committees to their union.

\subsection{Deferred evaluation claims}\label{sec:deferred-claims}

Twice during verification, a large \emph{fixed} multilinear must be evaluated at a random point:
\begin{itemize}
\item \textbf{the bytecode}: during leaf decomposition (\S\ref{sec:e2e-unrolled}, \emph{Bus}), the verifier evaluates the bytecode multilinear of \S\ref{sec:e2e-bc}, in $\kbc+4$ variables, in one pass over the whole program;
\item \textbf{the flock matrices}: the \texttt{BLAKE2S} relation is proven by flock (\S\ref{sec:tab-blake2s}, Annex~\ref{annex:flock}), whose lincheck ends by evaluating $\widetilde A_0$ and $\widetilde B_0$, the per-block R1CS matrices: $28$-variable multilinears with $\approx 89\mathrm{M}$ nonzero Boolean entries between them.
\end{itemize}
Either evaluation is one pass over a fixed object. This takes milliseconds natively, but in-circuit it would dwarf the entire remaining verifier, in cycles and in bytecode size alike.

The aggregation guest therefore does not evaluate these fixed polynomials. It verifies every surrounding reduction and exports only the resulting points and claimed values.

There are three fixed polynomials: the bytecode and the two matrices $A_0,B_0$. Each child contributes \emph{two} claims on each, the one it deferred and the fresh one that verifying it raises, so a node batches $2n_{\mathrm{rec}}$ per polynomial and a statement stays one size however deep the tree. A fresh Fiat--Shamir transcript binds all sub-statements and claim data, then batches them with two sumchecks: one for bytecode and one for the matrices, which share their $28$ variables. The aggregation and statement transcripts use distinct domain labels; each binds $n_{\mathrm{rec}}$ before any variable-length sequence, so concatenations cannot be confused. The guest reduces the batch to one bytecode evaluation and a common-point evaluation of each matrix, and those three enter its own statement, where its parent finds them.

A leaf has nothing to batch and defers the all-zeros point, where each polynomial is its table's first entry; the values ride an unchecked hint, since the verifier recomputes them from the point, so a lie changes the statement rather than the claim, and the same argument covers a carried claim. \texttt{AggregateSignature::verify} is the only acceptance path, and only the root pays those evaluations: every other node's claims are discharged by its parent's batch.

The matrix sumcheck runs over the $28$-variable polynomial
\[
  P(x) \;=\; A_0(x)\,W_A(x) \;+\; B_0(x)\,W_B(x),
  \qquad
  W_M(x) \;=\; \textstyle\sum_t \lambda_{M,t}\,\eq(r_t,x),
\]
and ends in the reduced claims $\widetilde A_0(r^\ast)$ and $\widetilde B_0(r^\ast)$, at a common point. The in-circuit cost is that of the sumcheck verifier: $28$ cheap rounds. A carried claim enters with the plain weight above. A fresh one enters with the weight of \eqref{flock:eq:matrix-form}, $e_{\mathrm{row}}\otimes w_{\mathrm{col}}$ on the row and column halves, at relative weight $\lcalpha$ of $B_0$ against $A_0$; its four tensor factors are evaluated at $r^\ast$ from the child's $\zskip$, $\rho_{\mathrm{in}}$, $s$ and $\rho'_{\mathrm{in}}$, all of which ride the aggregation statement. The prover never forms the $2^{28}$-sized boolean hypercube, and never reads a matrix entry either: both contractions it needs are circuit walks (\S\ref{flock:matrices}). The first $14$ rounds (the row variables) enter through the column contraction $m_t[i]=\sum_{j\in\mathrm{row}_i}\eq(r_t^{\mathrm{col}},j)$, which is the forward walk at the claim's column weights, keeping each row's value rather than contracting it; the last $14$ contract each matrix along its rows at the row point just fixed, which is the backward walk. Each is $O(\text{circuit})$, so the whole batch costs one forward walk per claim and two backward walks, plus $O(t\cdot2^{14})$ multiplications on dense tables of size $2^{14}$, against the $O(2^{28})$ field operations and memory of a dense sumcheck and the $\approx 89\mathrm{M}$ nonzeros a sparse one would iterate. The bytecode sumcheck is the same, only smaller, dense, and skipping the seven of sixteen stacking slots that are structurally zero.
